% Part: first-order-logic
% Chapter: syntax
% Section: intro-syntax

\documentclass[../../../include/open-logic-section]{subfiles}

\begin{document}

\olfileid[pt]{fol}{syn}{itx}

\olsection{Introdução}

Para desenvolver a teoria e a metateoria da lógica de primeira ordem,
devemos primeiro definir a sintaxe e a semântica de suas expressões.  As
expressões da lógica de primeira ordem são termos e !!{formula}s.  Termos
são formados a partir de !!{variable}s, !!{constant}s e !!{function}s.
!!^{formula}s, por sua vez, são formadas a partir de !!{predicate}s juntamente com
termos (estes formam as menores !!{formula}s, as ``atômicas''), e então,
a partir de !!{formula}s atômicas, podemos formar outras mais complexas usando conectivos
lógicos e quantificadores.  Há muitas maneiras diferentes de estabelecer
as regras de formação; damos apenas uma possível. Outros sistemas
escolherão símbolos diferentes, selecionarão conjuntos diferentes de
conectivos como primitivos, usarão parênteses de forma diferente (ou nem mesmo
os usarão, como no caso da chamada notação polonesa).  O que todas as
abordagens têm em comum, no entanto, é que as regras de formação definem
o conjunto de termos e !!{formula}s \emph{indutivamente}. Se feito
corretamente, cada expressão pode resultar essencialmente de uma única maneira
de acordo com as regras de formação.  A definição indutiva que resulta
em expressões que são \emph{de leitura única} significa que podemos dar
significados a essas expressões usando o mesmo método---a definição
indutiva.

\end{document}
